\documentclass[11pt, a4paper]{article}

% Gemeinsame Style-Definitionen (TEXINPUTS muss templates/ enthalten — siehe scripts/build_tex.py)
\input{bfw-style.tex}

\title{\vspace*{-1.5cm}\Huge\bfseries\color{bfwPrimaryDark}Session~2: Rechte, Pflichten \& Arbeitsrecht\\[0.4em]
  \large\color{bfwInkSoft}\textnormal{Von den BBiG-Pflichten bis zur Berufs- und Lebensplanung}}
\author{\textsf{IT Lernfeld 1 --- Das Unternehmen und die eigene Rolle im Betrieb beschreiben}}
\date{}

\begin{document}
\maketitle
\thispagestyle{empty}

\vspace{1em}
\begin{merksatz}[title={\bfseries Worum geht es in dieser Session?}]
Als Auszubildende im IT-Beruf stehen Sie in einem rechtlich geregelten Ausbildungsverhältnis.
Dieses Handout klärt, welche gegenseitigen Rechte und Pflichten das
\emph{Berufsbildungsgesetz (BBiG)} für Auszubildende und Ausbildende festlegt, welche
Schutzvorschriften das \emph{Jugendarbeitsschutzgesetz (JArbSchG)} für unter 18-Jährige
enthält und wie \emph{Betriebsrat} sowie \emph{Jugend- und Auszubildendenvertretung (JAV)}
die Interessen der Beschäftigten vertreten. Abschließend erhalten Sie einen Überblick über
das soziale Sicherungssystem und Instrumente zur persönlichen Berufs- und Lebensplanung.
\end{merksatz}

\tableofcontents
\vspace{1em}
\hrule
\vspace{1em}

\section{Rechte und Pflichten aus dem Ausbildungsverhältnis (§§~13–17 BBiG)}

Das Ausbildungsverhältnis ist ein gegenseitiges Rechtsverhältnis: Sowohl Ausbildende als auch
Auszubildende tragen Pflichten — und haben im Gegenzug Rechte. Die wichtigsten Regelungen
finden sich in den \textbf{§§~13 bis 17 BBiG}.

\subsection{Pflichten der Auszubildenden (§~13 BBiG)}

\begin{center}
\begin{tabular}{p{4cm} p{9cm}}
\toprule
\textbf{Pflicht} & \textbf{Erläuterung}\\
\midrule
Lernpflicht & Bemühen um Erwerb der erforderlichen Fertigkeiten, Kenntnisse und Erfahrungen\\
Teilnahmepflicht & Teilnahme an Ausbildungsmaßnahmen und Berufsschulunterricht (auch aktive Mitarbeit)\\
Weisungsgebundenheit & Befolgen der Anweisungen weisungsberechtigter Personen\\
Einhaltung der Ordnung & Beachten betrieblicher Ordnungsvorschriften\\
Ausbildungsnachweis & Schriftliche oder digitale Führung des Berichtshefts und regelmäßiges Vorlegen\\
Pfleglichkeit & Sorgfältige Behandlung der bereitgestellten Ausbildungsmittel\\
Benachrichtigungspflicht & Bei Fernbleiben: unverzügliche Nachricht mit Angabe von Gründen; bei Krankheit spätestens am dritten Tag ärztliche Bescheinigung\\
Geheimhaltungspflicht & Stillschweigen über Betriebs- und Geschäftsgeheimnisse während der gesamten Ausbildungszeit\\
\bottomrule
\end{tabular}
\end{center}

\subsection{Pflichten des Ausbildenden (§~14 BBiG)}

\begin{center}
\begin{tabular}{p{4cm} p{9cm}}
\toprule
\textbf{Pflicht} & \textbf{Erläuterung}\\
\midrule
Ausbildungspflicht & Planmäßige Vermittlung der zum Ausbildungsziel erforderlichen Fertigkeiten und Kenntnisse\\
Freistellungspflicht & Freistellung für Berufsschulunterricht, Prüfungen und Ausbildungsmaßnahmen (§§~14, 15 BBiG)\\
Benennung & Bekanntmachen der weisungsberechtigten Personen im Betrieb\\
Ausbildungsmittel & Kostenlose Bereitstellung aller für Ausbildung und Prüfung erforderlichen Mittel\\
Nachweiskontrolle & Aushändigen von Ausbildungsnachweisen und Überwachen ihrer ordnungsgemäßen Führung\\
Zweckbindung & Übertragung ausschließlich ausbildungszweckdienlicher und körperlich angemessener Aufgaben\\
Aufsichtspflicht & Beaufsichtigung minderjähriger Auszubildender während der betrieblichen Ausbildung\\
Zeugnispflicht & Schriftliches Zeugnis bei Beendigung der Ausbildung (§~16 BBiG)\\
\bottomrule
\end{tabular}
\end{center}

\begin{merksatz}[title={\bfseries Wichtig: Ausbildungszeugnis nach \S~16 BBiG}]
Das Zeugnis muss \textbf{Art, Dauer und Ziel} der Berufsausbildung sowie die erworbenen
\textbf{beruflichen Fertigkeiten, Kenntnisse und Fähigkeiten} enthalten. Angaben über
\textbf{Verhalten und Leistung} werden nur auf ausdrücklichen Wunsch des Auszubildenden
aufgenommen.
\end{merksatz}

\section{Ausbildungsvertrag, Vergütung, Probezeit und Kündigung}

\subsection{Mindestinhalte des Ausbildungsvertrags (§~11 BBiG)}

Vor Ausbildungsbeginn muss ein schriftlicher \textbf{Berufsausbildungsvertrag} abgeschlossen
werden. Er ist von allen Beteiligten zu unterzeichnen und bei der IHK einzureichen.
§~11 BBiG schreibt \textbf{zehn Pflichtinhalte} vor:

\begin{enumerate}
  \item Art, sachliche und zeitliche Gliederung sowie Ziel der Berufsausbildung
  \item Beginn und Dauer der Berufsausbildung
  \item Ausbildungsmaßnahmen außerhalb der Ausbildungsstätte
  \item Dauer der regelmäßigen täglichen Ausbildungszeit
  \item Dauer der Probezeit
  \item Zahlung und Höhe der Vergütung
  \item Dauer des Urlaubs
  \item Voraussetzungen, unter denen der Vertrag gekündigt werden kann
  \item Hinweis auf anzuwendende Tarifverträge oder Betriebsvereinbarungen
  \item Form des Ausbildungsnachweises (schriftlich oder digital)
\end{enumerate}

Vereinbarungen, die den Vorschriften des BBiG widersprechen oder dem Zweck der
Berufsausbildung zuwiderlaufen, sind nichtig --- unabhängig davon, ob beide Seiten
zugestimmt haben.

\subsection{Ausbildungsvergütung (§§~17, 18 BBiG)}

Die Ausbildungsvergütung muss \textbf{angemessen} sein und \textbf{mindestens jährlich
steigen}. Seit der Einführung einer gesetzlichen Mindestausbildungsvergütung (§~17 BBiG) ist
eine Untergrenze festgelegt. Bei Krankheit wird die Vergütung nach dem
Entgeltfortzahlungsgesetz (EntgFG) bis zu sechs Wochen weitergezahlt, danach zahlt die
Krankenkasse Krankengeld. In IT-Berufen liegen die Vergütungen in der Regel deutlich über
der gesetzlichen Mindestgrenze.

\begin{center}
\begin{tabular}{p{5cm} p{8cm}}
\toprule
\textbf{Abzug} & \textbf{Erläuterung}\\
\midrule
Lohn- und Kirchensteuer & Abhängig von Bruttobetrag und Steuerklasse; Arbeitgeber überweist ans Finanzamt\\
Sozialversicherungsbeiträge & Kranken-, Pflege-, Renten- und Arbeitslosenversicherung; je zur Hälfte von AG und AN getragen\\
Unfallversicherung & Trägt der Arbeitgeber allein (Berufsgenossenschaft)\\
\bottomrule
\end{tabular}
\end{center}

\subsection{Probezeit (§~20 BBiG)}

Jedes Ausbildungsverhältnis beginnt mit einer \textbf{Probezeit} von mindestens einem und
höchstens vier Monaten. Während der Probezeit kann \emph{jede Seite} ohne Frist und ohne
Angabe von Gründen kündigen.

\subsection{Kündigung nach der Probezeit (§~22 BBiG)}

\begin{center}
\begin{tabular}{p{5cm} p{8cm}}
\toprule
\textbf{Partei} & \textbf{Kündigungsmöglichkeit nach der Probezeit}\\
\midrule
Ausbildende:r & Nur außerordentliche fristlose Kündigung aus \emph{wichtigem Grund} möglich; keine ordentliche Kündigung\\
Auszubildende:r & Fristgerecht mit \textbf{vier Wochen} Kündigungsfrist, wenn: (1) Aufgabe der Berufsausbildung oder (2) Berufswechsel angestrebt wird; außerdem außerordentlich fristlos aus wichtigem Grund\\
\bottomrule
\end{tabular}
\end{center}

\begin{merksatz}
Eine fristlose Kündigung aus wichtigem Grund muss \textbf{innerhalb von zwei Wochen} nach
Kenntnis des Kündigungsgrundes ausgesprochen werden, sonst ist sie unwirksam (§~22
Abs.~4 BBiG).
\end{merksatz}

\section{Jugendarbeitsschutzgesetz (JArbSchG)}

Das \textbf{Jugendarbeitsschutzgesetz} gilt für alle Personen unter 18 Jahren in
Berufsausbildung oder Dauerbeschäftigung von mehr als zwei Monaten außerhalb des
Familienhaushalts.

\begin{center}
\begin{tabular}{p{4cm} p{9cm}}
\toprule
\textbf{Regelung} & \textbf{Inhalt}\\
\midrule
Tägliche Arbeitszeit & Höchstens 8 Stunden; in Ausnahmen 8{,}5 Stunden (§~8 JArbSchG)\\
Wöchentliche Arbeitszeit & Höchstens 40 Stunden, fünf Arbeitstage\\
Beschäftigungszeit & Nur von 6 bis 20 Uhr; in Mehrschichtbetrieben bis 23 Uhr\\
Ruhepausen & Bei mehr als 4{,}5–6 Stunden: mind. 30 Minuten; bei mehr als 6 Stunden: mind. 60 Minuten\\
Ruhezeit & Mindestens 12 Stunden Freizeit nach Ende der Arbeitszeit\\
Urlaubsanspruch & Unter 16 Jahre: 30 Tage; unter 17 Jahre: 27 Tage; unter 18 Jahre: 25 Tage\\
Mehrarbeit & Muss innerhalb der folgenden drei Wochen durch Freizeit ausgeglichen werden\\
Nachtarbeit & In IT-Betrieben grundsätzlich nicht erlaubt\\
Prüfungsfreistellung & Freizustellen am Tag vor der schriftlichen Prüfung\\
Schulfreistellung & Ganztägig freizustellen, wenn Unterricht an einem Tag mehr als fünf Stunden umfasst (einmal/Woche)\\
Erstuntersuchung & Vor Beginn der Ausbildung; nach einem Jahr Nachuntersuchung (§§~32 ff.~JArbSchG)\\
\bottomrule
\end{tabular}
\end{center}

\section{Arbeitsrecht und Mitbestimmung im Betrieb}

\subsection{Wesentliche Arbeitsgesetze}

Der Staat hat eine Vielzahl von Gesetzen erlassen, um die Interessen von Arbeitnehmern zu
schützen. Grundsätzlich gilt die \textbf{Vertragsfreiheit}: Arbeitsverträge können individuell
gestaltet werden. Wo jedoch Gesetze, Tarifverträge oder Betriebsvereinbarungen bestehen,
gilt das \textbf{Günstigkeitsprinzip}: Stets hat die für den Arbeitnehmer günstigere Regelung
Vorrang.

\begin{center}
\begin{tabular}{p{4cm} p{9cm}}
\toprule
\textbf{Gesetz} & \textbf{Kerninhalt}\\
\midrule
ArbZG & Begrenzt die höchstzulässige tägliche Arbeitszeit für Erwachsene (8 Std., bis 10 Std. möglich); Mindestruhezeiten\\
KSchG / §~622 BGB & Kündigungsschutz; ab mehr als zehn Vollbeschäftigten und sechs Monaten Betriebszugehörigkeit\\
AGG & Verbot der Diskriminierung aufgrund Geschlecht, Herkunft, Alter, Religion, Behinderung, sexueller Identität\\
MiLoG & Gesetzliche Mindestlohngrenze; Mindestausbildungsvergütung nach § 17 BBiG\\
BurlG & Mindesturlaub 24 Werktage (bei 5 Arbeitstagen/Woche: 20 Tage)\\
MuSchG / BEEG & Mutterschutz (Schutzfrist 6 Wochen vor / 8 Wochen nach Entbindung); Elternzeit bis zum 3. Lebensjahr\\
EntgFG & Entgeltfortzahlung bei Erkrankung für bis zu 6 Wochen\\
\bottomrule
\end{tabular}
\end{center}

\subsection{Betriebsrat (BetrVG)}

Das \textbf{Betriebsverfassungsgesetz (BetrVG)} regelt die Zusammenarbeit zwischen
Arbeitgeber und Betriebsrat. In jedem Betrieb mit mindestens \textbf{fünf wahlberechtigten
Beschäftigten} (ab 18 Jahre) kann ein Betriebsrat gewählt werden. Die
\textbf{Amtszeit beträgt vier Jahre}. Der Betriebsrat:

\begin{itemize}
  \item verhandelt mit dem Arbeitgeber über Arbeitszeit, Urlaub, Lohn und Arbeitsbedingungen,
  \item hat bei bestimmten Maßnahmen (z.\,B.\ Kündigungen, Neueinstellungen) ein
    Mitbestimmungsrecht,
  \item beruft mindestens einmal pro Quartal eine Betriebsversammlung ein (§~43 BetrVG),
  \item kann Betriebsvereinbarungen mit dem Arbeitgeber abschließen.
\end{itemize}

\begin{merksatz}
Betriebsratsmitglieder genießen noch bis \textbf{12 Monate} nach Ende ihrer Amtszeit
besonderen Kündigungsschutz.
\end{merksatz}

\subsection{Jugend- und Auszubildendenvertretung (JAV)}

Die \textbf{JAV} ist die gewählte Interessenvertretung aller Jugendlichen und
Auszubildenden im Betrieb. Sie setzt voraus, dass \textbf{mindestens fünf wahlberechtigte
Jugendliche oder Auszubildende} im Betrieb beschäftigt sind.

Wahlberechtigt sind:
\begin{itemize}
  \item alle Beschäftigten unter 18 Jahren und
  \item alle Auszubildenden bis zum vollendeten 25. Lebensjahr.
\end{itemize}

Die JAV arbeitet eng mit dem Betriebsrat zusammen, überwacht die Einhaltung von BBiG,
JArbSchG und Betriebsvereinbarungen und setzt sich für die Übernahme der Auszubildenden
nach der Ausbildung ein. Gemäß §~78a BetrVG haben JAV-Mitglieder bei Ausbildungsende
ein \textbf{Recht auf Übernahme in ein unbefristetes Arbeitsverhältnis}.

\subsection{Tarifverträge}

Tarifverträge werden zwischen \textbf{Gewerkschaften} und \textbf{Arbeitgeberverbänden}
(oder einzelnen Unternehmen als Tarifparteien) abgeschlossen und regeln Arbeits- und
Einkommensbedingungen. Sie gelten grundsätzlich für tarifgebundene Mitglieder. Das
Bundesministerium für Arbeit und Soziales kann Tarifverträge für \textbf{allgemeinverbindlich}
erklären --- dann gelten sie für alle Betriebe des Geltungsbereichs (§~5 TVG).

\section{Grundzüge des sozialen Sicherungssystems}

Das soziale Sicherungssystem der Bundesrepublik Deutschland beruht auf den Prinzipien der
\textbf{sozialen Marktwirtschaft}, der \textbf{Solidarität} (Lasten auf viele Schultern
verteilen) und der \textbf{Subsidiarität} (Eigenverantwortung vor staatlichem Eingriff).

\begin{center}
\begin{tabular}{p{4cm} p{9cm}}
\toprule
\textbf{Säule} & \textbf{Träger und Funktion}\\
\midrule
Krankenversicherung & Gesetzliche Krankenkassen; deckt Arzt- und Behandlungskosten ab\\
Pflegeversicherung & Gesetzliche Krankenkassen; finanzielle Unterstützung bei Pflegebedürftigkeit\\
Rentenversicherung & Deutsche Rentenversicherung; monatliche Altersrente ab Rentenalter\\
Arbeitslosenversicherung & Bundesagentur für Arbeit; Lohnersatz bei Jobverlust (Arbeitslosengeld I)\\
Unfallversicherung & Berufsgenossenschaften; Absicherung bei Arbeitsunfällen --- Kosten trägt allein der Arbeitgeber\\
\bottomrule
\end{tabular}
\end{center}

Beiträge zur Kranken-, Pflege-, Renten- und Arbeitslosenversicherung werden jeweils zur
Hälfte von Arbeitgeber und Arbeitnehmer getragen und vom Arbeitgeber direkt überwiesen.
Zusätzlich können private Versicherungen (Haftpflicht, Berufsunfähigkeit, Lebensversicherung)
sinnvoll sein, um Risiken jenseits der gesetzlichen Grundsicherung abzudecken.

\section{Berufs- und Lebensplanung}

\begin{merksatz}[title={\bfseries Lebenslanges Lernen als Schlüssel}]
In der IT-Branche verändern sich Technologien und Anforderungen besonders schnell.
Wer erfolgreich bleiben will, muss die eigene Karriere und Weiterbildung aktiv gestalten
--- durch Fortbildung, Flexibilität und vorausschauende Lebensplanung.
\end{merksatz}

\subsection{Fortbildungsarten nach BBiG}

Das BBiG unterscheidet folgende Fortbildungsarten:
\begin{itemize}
  \item \textbf{Erhaltungsfortbildung}: vorhandene Kenntnisse auffrischen
  \item \textbf{Anpassungsfortbildung}: an neue Anforderungen anpassen
  \item \textbf{Aufstiegsfortbildung}: für höhere Positionen qualifizieren
  \item \textbf{Innerbetriebliche Fortbildung}: betriebsintern organisiert
  \item \textbf{Überbetriebliche Fortbildung}: in externen Einrichtungen
\end{itemize}

\subsection{Karriere- und Lebensplanung}

\begin{center}
\begin{tabular}{p{4cm} p{9cm}}
\toprule
\textbf{Instrument} & \textbf{Erläuterung}\\
\midrule
Potenzialanalyse & Erfassung eigener Stärken und Schwächen als Planungsgrundlage\\
Karriereplanung & Vorausschauende Gestaltung des Berufswegs; Prüfung von Erfolgsfaktoren\\
Europass & Europäisches Karriereportal zur mehrsprachigen Dokumentation von Qualifikationen\\
Erasmus+ & EU-Förderprogramm für Auslandsaufenthalte auch während der Berufsausbildung\\
Job-Rotation & Wechsel des Arbeitsplatzes zur Vertiefung von Wissen und Vermeidung von Monotonie\\
Job-Enrichment & Erhöhung von Verantwortung und Entscheidungsspielraum zur Motivationsförderung\\
Private Altersvorsorge & Ergänzung zur gesetzlichen Rente durch Immobilien, Kapital, Betriebs- oder Privatrenten\\
\bottomrule
\end{tabular}
\end{center}

\section{Zusammenfassung}

\begin{enumerate}
  \item \textbf{§§~13, 14 BBiG} regeln die gegenseitigen Pflichten: Auszubildende (Lernpflicht, Geheimhaltung, Berichtsheft) und Ausbildende (Ausbildungspflicht, Freistellung, Mittelbereitstellung).
  \item \textbf{§~11 BBiG} schreibt zehn Pflichtinhalte des Ausbildungsvertrags vor; Verstöße gegen das BBiG machen Vereinbarungen nichtig.
  \item \textbf{§§~17, 18 BBiG}: Angemessene, jährlich steigende Vergütung.
  \item \textbf{§~20 BBiG}: Probezeit 1–4 Monate.
  \item \textbf{§~22 BBiG}: Kündigung nach der Probezeit durch Ausbildende nur außerordentlich fristlos.
  \item \textbf{JArbSchG}: Schutz unter 18-Jähriger --- tägliche Arbeitszeit max.\ 8 Stunden, gestaffelter Urlaub, Erstuntersuchung, Nachtarbeitsverbot, Ausgleich von Mehrarbeit.
  \item \textbf{Betriebsrat (BetrVG)}: Ab fünf Beschäftigten wählbar; vier Jahre Amtszeit; Mitbestimmungsrechte bei Kündigungen und Neueinstellungen.
  \item \textbf{JAV (§~70 BetrVG)}: Vertretung der Jugendlichen und Auszubildenden; überwacht BBiG- und JArbSchG-Einhaltung.
  \item Das soziale Sicherungssystem beruht auf fünf \textbf{Säulen}: Kranken-, Pflege-, Renten-, Arbeitslosen- und Unfallversicherung.
  \item Persönliche \textbf{Berufs- und Lebensplanung} umfasst Potenzialanalyse, Fortbildung, Mobilität (Europass, Erasmus+) und private Altersvorsorge.
\end{enumerate}

\newpage

\section*{Glossar}
\addcontentsline{toc}{section}{Glossar}

\glossareintrag{BBiG}{Berufsbildungsgesetz; regelt Rechte und Pflichten in der Berufsausbildung, Ausbildungsvertrag, Vergütung, Probezeit und Kündigung.}

\glossareintrag{JArbSchG}{Jugendarbeitsschutzgesetz; schützt Beschäftigte unter 18 Jahren durch Begrenzung der Arbeitszeit, Urlaubsregelungen und Untersuchungspflichten.}

\glossareintrag{Geheimhaltungspflicht}{Pflicht des Auszubildenden nach § 13 BBiG, über Betriebs- und Geschäftsgeheimnisse Stillschweigen zu wahren.}

\glossareintrag{Probezeit}{Eingangsphase des Ausbildungsverhältnisses von mindestens einem und höchstens vier Monaten; in dieser Zeit kann jede Seite ohne Frist kündigen (§ 20 BBiG).}

\glossareintrag{BetrVG}{Betriebsverfassungsgesetz; regelt die Mitbestimmung der Arbeitnehmer im Betrieb durch den Betriebsrat.}

\glossareintrag{Betriebsrat}{Gewählte Interessenvertretung der Arbeitnehmer; ab fünf wahlberechtigten Beschäftigten möglich; Amtszeit vier Jahre.}

\glossareintrag{JAV}{Jugend- und Auszubildendenvertretung; gewählte Interessenvertretung aller Jugendlichen und Auszubildenden im Betrieb; arbeitet mit dem Betriebsrat zusammen.}

\glossareintrag{Günstigkeitsprinzip}{Grundsatz im Arbeitsrecht: Stets gilt die für den Arbeitnehmer günstigere Regelung, unabhängig davon, ob sie aus Gesetz, Tarifvertrag oder Arbeitsvertrag stammt.}

\glossareintrag{Tarifvertrag}{Vereinbarung zwischen Gewerkschaft und Arbeitgeberverband über Arbeits- und Einkommensbedingungen; kann für allgemeinverbindlich erklärt werden.}

\glossareintrag{Subsidiarität}{Gesellschaftliches Prinzip, nach dem Aufgaben von der kleinsten geeigneten Einheit selbst gelöst werden sollen; Grundsatz der sozialen Marktwirtschaft.}

\glossareintrag{Europass}{Europäisches Karriereportal zur mehrsprachigen und datensicheren Dokumentation von Qualifikationen, Abschlüssen und Lern- sowie Arbeitserfahrungen.}

\end{document}
